\chapter{Literature Review}

Cite the following: International allocation and currency risk management, Currency overlay vs. joint optimization, Downside risk measures and Expected Shortfall, Instability of ES optimization and the role of regularization, Generalized objectives and decision-focused / end-to-end learning.

In this chapter, we review the relevant literature on multi-currency 
portfolio optimization. We discuss the contributions and 
methodologies in the field, highlighting how our work builds 
upon and extends existing research. 

The core of our analysis builds on two recent studies that 
apply regularization techniques to international asset allocation. 
Burkhardt and Ulrych (2023) introduce a joint optimization approach 
within a mean-variance framework. By using L1 and L2 regularization, they show that jointly optimizing assets and 
currencies leads to better out-of-sample performance than the standard 
industry practice of using a separate currency overlay. 

Extending this framework, Lucescu and Ulrych (2026) consider an 
Expected Shortfall optimization framework to better manage downside 
tail risk. To fix the numerical instability common in Expected Shortfall minimization, 
they apply similar regularization methods. Their findings reveal a 
dimensional trade-off. Joint optimization works best in small asset universes, 
but the separate currency overlay approach becomes more effective for larger 
universes due to its built-in dimensionality reduction. 

Reproducing and comparing these two baseline frameworks serves as the starting point 
for the extensions proposed in this thesis.


% =============================================================================
% STRUCTURE INSTRUCTIONS (delete before submission)
% Nothing above this line has been removed -- the prose on the two baseline
% papers is kept and should be reused, most of it inside block (4) below.
%
% WHY THIS CHAPTER NEEDS RESTRUCTURING
% The contribution in Chapter~\ref{chap:joint-vs-sequential} is an
% ESTIMATION-ERROR result: the joint program's in-sample advantage over the
% sequential overlay is largest precisely where it is least realizable,
% because the measures that produce the largest in-sample gap are estimated
% from the fewest observations. If this review covers only currency hedging
% and risk measures, that contribution has no home when the reader reaches it.
% The reader must already be standing in the estimation-error literature.
%
% PROPOSED BLOCKS, in this order:
%
% (1) MULTI-CURRENCY PORTFOLIO CONSTRUCTION
%     The return decomposition and the currency-exposure problem. Set up why
%     currency is a separate decision at all. Existing refs:
%     global_currency_hedging, CHEN2003133, safe_haven_currencies,
%     common_risk_factors_currency_markets,
%     carry_trades_and_global_foreign_exchange_volatility,
%     MenkhoffSarnoSchmelingSchrimpf2012, NeelyWellerUlrich2009.
%
% (2) RISK MEASURES BEYOND VARIANCE
%     Markowitz1952 as the origin; CoherentMeasuresOfRisk for the coherence
%     axioms; Rockafellar-Uryasev for the CVaR linear program (ADD -- not in
%     refs.bib yet and it is load-bearing for Chapter 3); Konno-Yamazaki for
%     MAD (ADD); Krokhmal for higher-moment coherent risk measures / HMCR_2
%     (ADD -- this is your SOES). goldberg2013minimizingshortfall fits here.
%     End the block by naming the property your results turn on: deviation
%     measures use every observation, tail measures use only the tail.
%
% (3) ESTIMATION ERROR IN PORTFOLIO OPTIMIZATION  <-- THE LOAD-BEARING BLOCK
%     This is where your contribution lands, so it deserves the most space
%     and currently has the least. Existing refs to build on: Michaud1989
%     (the "estimation-error maximiser"), Jorion1986BayesStein, LedoitWolf2004,
%     ciliberti2007feasibility and kondor2007noise (noise sensitivity of ES
%     optimization specifically -- these are the closest existing work to your
%     level-sweep result and should be discussed, not just cited).
%     ADD: DeMiguel, Garlappi & Uppal (2009) on 1/N versus optimized
%     portfolios, which is the canonical statement that in-sample optimality
%     need not survive; and DeMiguel et al. on norm-constrained portfolios,
%     which is the regularization-as-estimation-control argument.
%     The paragraph to end on: the literature establishes that optimized
%     portfolios overfit, and that the effect scales with dimension. What it
%     does not establish is that the SIZE of an in-sample advantage predicts
%     its realizability NEGATIVELY across risk measures. State that gap
%     explicitly -- it is your contribution's entry point.
%
% (4) HEDGING: OVERLAY VERSUS INTEGRATED
%     The two baseline papers, using the prose already written above.
%     SparseAndStableInternationalPortfolioOptimization (Burkhardt & Ulrych,
%     mean-variance, joint beats overlay out of sample with regularization)
%     and RegularizedMultiCurrencyExpectedShortfallPortfolios (Lucescu &
%     Ulrych, ES, the dimensional trade-off).
%     IMPORTANT: state precisely what each paper claims and on which universe
%     sizes, because your results agree with the second and qualify the first.
%     Then state what this thesis adds: five risk measures rather than one or
%     two, a decomposition of the in-sample advantage into a structural part
%     and an estimation part, and the finding that the in-sample gap is a
%     measure of overfitting capacity rather than of economic value.
%
% REFERENCES TO ADD TO refs.bib (checked -- none are present):
%   Rockafellar & Uryasev (2000), Optimization of Conditional Value-at-Risk
%   Konno & Yamazaki (1991), Mean-Absolute Deviation Portfolio Optimization
%   Krokhmal (2007), Higher Moment Coherent Risk Measures
%   DeMiguel, Garlappi & Uppal (2009), Optimal Versus Naive Diversification
%   DeMiguel, Garlappi, Nogales & Uppal (2009), Norm-Constrained Portfolios
%
% The line at the top of this file listing topics to cite is superseded by the
% blocks above; delete it once the blocks are written.
% =============================================================================
